\section{Specification history and the placebo test}\label{app:history}

We report the specifications that \emph{failed} the built-in falsification test, because
the failures are informative and because the final specification was chosen by them, not
by the results we hoped for. All superseded outputs are kept in
\texttt{results/json/superseded/}.

\paragraph{Version 1.} This version used news kernels with time scales up to 30 minutes and
Erlang order 3, so that mean lags reached 90 minutes. It had a constant baseline in each window
and time-of-day effects in UTC, and it allowed tied events: a trade that crossed several $\delta$
levels produced several events with the same millisecond timestamp. The single placebo type
absorbed \vOnePlacebo{} extra events per pseudo-release, which is comparable to real releases, and
the residuals of the crypto series had Kolmogorov--Smirnov statistics up to \vOneKSmax. Two causes
explain this. News kernels longer than the post-release window absorbed slow trends within the
window, and ties, which made up 55\% of bitcoin events, violate the assumption that events occur
one at a time.

\paragraph{Version 2.} This version kept one event per timestamp, with the number of crossed
levels as a mark. It restricted news kernels to time scales between 0.1 and 300\,s, used a linear
baseline within each window and modelled seasonality in New York time. The placebo mass fell to
\vTwoPlacebo{} and the residual fit improved by an order of magnitude, but the placebo effect was
still significant under Wald standard errors. Splitting the placebo by time of day showed placebo
mass at every release time, including 14:00\,ET on days without an FOMC meeting. Unmodelled 08:30
releases therefore could not be the whole explanation.

\paragraph{Version 3 (final).} This version made two further changes. First, placebo days must also
avoid sixteen second-tier US releases, such as PCE, durable goods, trade, housing, JOLTS and
consumer sentiment. Second, the news kernels carry an $\ell_1$ penalty, weighted by exposure, whose
level is chosen by the likelihood of held-out windows (Table~\ref{tab:l1}). Without the penalty,
each of the several hundred nonnegative weights has a positive expected value even when there is
no effect, so release and placebo masses are inflated equally. The held-out likelihood peaks at
$\ell_1=\lOneSelected$, which is also where release and placebo masses separate.

\begin{table}[!htbp]\centering\small
\caption{Choice of the $\ell_1$ penalty on the news kernels. Held-out log-likelihood on the
2024 windows relative to the best level (bold), and the mean news-kernel mass of release and
placebo windows.}\label{tab:l1}
\begin{tabular}{cccc}
\toprule
$\ell_1$ level & held-out LL (rel.) & mean release mass & mean placebo mass \\
\midrule
0.2 & {$-$93.7} & 10.79 & 6.23 \\
0.4 & {$-$27.1} & 5.20 & 2.89 \\
0.8 & \textbf{0.0} & 1.81 & 1.03 \\
1.6 & {$-$20.5} & 0.44 & 0.37 \\
\bottomrule
\end{tabular}
\end{table}

\section{EPT ablations}
\paragraph{Hawkes backbone.} Without the explicit linear Hawkes state, EPT reached
\eptNoBackbone{} nats per event on LOBSTER AAPL (three seeds, 100 epochs), which is below the
classical MSX estimator. The network had to rediscover mark-specific linear excitation through
its gated jumps, and it did not manage to within the epoch budget. The backbone is a residual
path that makes the model exactly MSX when the neural read-outs are zero. With it, EPT is better
than every neural baseline on all five stocks and better than MSX on three of them
(Table~\ref{tab:track-c}).

\paragraph{Renewal channel and encoder.} Adding the mixture-of-Erlang renewal channel (EPT-X)
improves EPT on every LOBSTER stock and on every public dataset where both were run. On
validation, sharper Erlang atoms (orders up to 1024) help where the gap distribution has sharp
peaks: by 1.1 nats per event on Amazon, whose gaps are concentrated in a narrow band, and by 0.05
on INTC. Replacing the GRU by a causal Transformer encoder \citep{vaswani2017attention} lowers the
validation likelihood on all five datasets we tested, and neither a three-layer residual encoder
nor a 128-unit encoder improves any dataset (Section~\ref{sec:results}).

\clearpage
\section{Track A: all metrics}
\begin{center}\footnotesize
\captionof{table}{Synthetic recovery, all metrics (mean $\pm$ s.d.\ over five seeds):
relative error of the branching matrix and of its spectral radius, kernel $L^1$ error and held-out
deficit dLL ($10^{-3}$ nats/event), all lower is better, and edge AUC (higher is better). AUC needs
absent true edges, so it is defined only for S4 and S5 (n/a for the complete graphs S1--S3).
``--'': not produced by the method; NPHC estimates only the branching matrix. Missing rows are fits
that did not converge. Bold: best; underlined: second best within each scenario; shaded: ours.}\label{tab:track-a-full}\smallskip
\scriptsize\setlength{\tabcolsep}{5pt}\renewcommand{\arraystretch}{1.0}
\begin{tabular}{lrrrrr}
\toprule
& $G$ error & $\rho$ error & kernel $L^1$ & dLL ($10^{-3}$) & edge AUC \\
\midrule
\multicolumn{6}{l}{\textit{S1: exponential}} \\[1pt]
\quad Exp.\ kernel MLE & \textbf{0.018}{\scriptsize\,$\pm$0.007} & \textbf{0.005}{\scriptsize\,$\pm$0.003} & \textbf{0.018}{\scriptsize\,$\pm$0.008} & \textbf{0.05}{\scriptsize\,$\pm$0.05} & n/a \\
\quad Sum-of-exp.\ MLE & 0.065{\scriptsize\,$\pm$0.022} & 0.030{\scriptsize\,$\pm$0.013} & 0.165{\scriptsize\,$\pm$0.023} & 0.92{\scriptsize\,$\pm$0.13} & n/a \\
\quad ADM4 & \underline{0.018}{\scriptsize\,$\pm$0.007} & \underline{0.005}{\scriptsize\,$\pm$0.003} & \underline{0.018}{\scriptsize\,$\pm$0.008} & \underline{0.05}{\scriptsize\,$\pm$0.05} & n/a \\
\quad Hawkes EM (nonparam.) & 0.035{\scriptsize\,$\pm$0.015} & 0.017{\scriptsize\,$\pm$0.007} & 0.173{\scriptsize\,$\pm$0.010} & 1.02{\scriptsize\,$\pm$0.07} & n/a \\
\quad NPHC (cumulants) & 0.247{\scriptsize\,$\pm$0.087} & 0.032{\scriptsize\,$\pm$0.024} & -- & -- & n/a \\
\quad Conditional law & 0.120{\scriptsize\,$\pm$0.054} & 0.017{\scriptsize\,$\pm$0.016} & 0.193{\scriptsize\,$\pm$0.022} & 2.84{\scriptsize\,$\pm$1.03} & n/a \\
\rowcolor{oursbg} \quad MSX-auto & 0.054{\scriptsize\,$\pm$0.032} & 0.020{\scriptsize\,$\pm$0.018} & 0.093{\scriptsize\,$\pm$0.035} & 0.24{\scriptsize\,$\pm$0.12} & n/a \\
\rowcolor{oursbg} \quad MSX ($R=2$) & 0.056{\scriptsize\,$\pm$0.032} & 0.019{\scriptsize\,$\pm$0.018} & 0.092{\scriptsize\,$\pm$0.035} & 0.23{\scriptsize\,$\pm$0.12} & n/a \\
\rowcolor{oursbg} \quad MSX ($R=4$) & 0.056{\scriptsize\,$\pm$0.035} & 0.023{\scriptsize\,$\pm$0.018} & 0.105{\scriptsize\,$\pm$0.034} & 0.28{\scriptsize\,$\pm$0.15} & n/a \\
\rowcolor{oursbg} \quad MSX, exponential only & 0.055{\scriptsize\,$\pm$0.025} & 0.026{\scriptsize\,$\pm$0.015} & 0.157{\scriptsize\,$\pm$0.022} & 0.91{\scriptsize\,$\pm$0.13} & n/a \\
\addlinespace[3pt]\multicolumn{6}{l}{\textit{S2: power law}} \\[1pt]
\quad Exp.\ kernel MLE & 0.299{\scriptsize\,$\pm$0.011} & 0.219{\scriptsize\,$\pm$0.007} & 0.465{\scriptsize\,$\pm$0.003} & 26.05{\scriptsize\,$\pm$0.50} & n/a \\
\quad Sum-of-exp.\ MLE & \underline{0.057}{\scriptsize\,$\pm$0.012} & 0.016{\scriptsize\,$\pm$0.007} & \textbf{0.094}{\scriptsize\,$\pm$0.011} & \textbf{0.18}{\scriptsize\,$\pm$0.12} & n/a \\
\quad ADM4 & 0.299{\scriptsize\,$\pm$0.011} & 0.219{\scriptsize\,$\pm$0.007} & 0.465{\scriptsize\,$\pm$0.003} & 26.05{\scriptsize\,$\pm$0.50} & n/a \\
\quad Hawkes EM (nonparam.) & 0.057{\scriptsize\,$\pm$0.014} & 0.037{\scriptsize\,$\pm$0.008} & 0.188{\scriptsize\,$\pm$0.006} & 1.73{\scriptsize\,$\pm$0.21} & n/a \\
\quad NPHC (cumulants) & 0.153{\scriptsize\,$\pm$0.154} & 0.041{\scriptsize\,$\pm$0.047} & -- & -- & n/a \\
\quad Conditional law & 0.102{\scriptsize\,$\pm$0.067} & 0.029{\scriptsize\,$\pm$0.019} & 0.291{\scriptsize\,$\pm$0.028} & 4.91{\scriptsize\,$\pm$0.59} & n/a \\
\rowcolor{oursbg} \quad MSX-auto & 0.064{\scriptsize\,$\pm$0.029} & 0.037{\scriptsize\,$\pm$0.029} & 0.108{\scriptsize\,$\pm$0.022} & 0.33{\scriptsize\,$\pm$0.13} & n/a \\
\rowcolor{oursbg} \quad MSX ($R=2$) & 0.057{\scriptsize\,$\pm$0.021} & \underline{0.015}{\scriptsize\,$\pm$0.010} & 0.114{\scriptsize\,$\pm$0.015} & 0.32{\scriptsize\,$\pm$0.12} & n/a \\
\rowcolor{oursbg} \quad MSX ($R=4$) & 0.062{\scriptsize\,$\pm$0.014} & 0.016{\scriptsize\,$\pm$0.006} & 0.131{\scriptsize\,$\pm$0.017} & 0.45{\scriptsize\,$\pm$0.10} & n/a \\
\rowcolor{oursbg} \quad MSX, exponential only & \textbf{0.054}{\scriptsize\,$\pm$0.022} & \textbf{0.014}{\scriptsize\,$\pm$0.010} & \underline{0.097}{\scriptsize\,$\pm$0.011} & \underline{0.19}{\scriptsize\,$\pm$0.11} & n/a \\
\addlinespace[3pt]\multicolumn{6}{l}{\textit{S3: hump}} \\[1pt]
\quad Exp.\ kernel MLE & 0.166{\scriptsize\,$\pm$0.008} & 0.053{\scriptsize\,$\pm$0.004} & 0.801{\scriptsize\,$\pm$0.004} & 101.47{\scriptsize\,$\pm$1.81} & n/a \\
\quad Sum-of-exp.\ MLE & 0.038{\scriptsize\,$\pm$0.011} & 0.018{\scriptsize\,$\pm$0.005} & 0.419{\scriptsize\,$\pm$0.005} & 11.51{\scriptsize\,$\pm$0.46} & n/a \\
\quad ADM4 & 0.166{\scriptsize\,$\pm$0.008} & 0.053{\scriptsize\,$\pm$0.004} & 0.801{\scriptsize\,$\pm$0.004} & 101.47{\scriptsize\,$\pm$1.81} & n/a \\
\quad Hawkes EM (nonparam.) & \textbf{0.024}{\scriptsize\,$\pm$0.011} & \textbf{0.010}{\scriptsize\,$\pm$0.004} & 0.187{\scriptsize\,$\pm$0.008} & 2.24{\scriptsize\,$\pm$0.12} & n/a \\
\quad NPHC (cumulants) & 0.240{\scriptsize\,$\pm$0.096} & 0.045{\scriptsize\,$\pm$0.048} & -- & -- & n/a \\
\quad Conditional law & 0.149{\scriptsize\,$\pm$0.062} & 0.017{\scriptsize\,$\pm$0.010} & 0.319{\scriptsize\,$\pm$0.050} & 10.66{\scriptsize\,$\pm$2.13} & n/a \\
\rowcolor{oursbg} \quad MSX-auto & \underline{0.032}{\scriptsize\,$\pm$0.016} & 0.016{\scriptsize\,$\pm$0.007} & \underline{0.080}{\scriptsize\,$\pm$0.013} & \textbf{0.30}{\scriptsize\,$\pm$0.13} & n/a \\
\rowcolor{oursbg} \quad MSX ($R=2$) & 0.069{\scriptsize\,$\pm$0.009} & 0.037{\scriptsize\,$\pm$0.004} & 0.347{\scriptsize\,$\pm$0.007} & 5.90{\scriptsize\,$\pm$0.36} & n/a \\
\rowcolor{oursbg} \quad MSX ($R=4$) & 0.038{\scriptsize\,$\pm$0.010} & 0.017{\scriptsize\,$\pm$0.005} & \textbf{0.074}{\scriptsize\,$\pm$0.010} & \underline{0.31}{\scriptsize\,$\pm$0.15} & n/a \\
\rowcolor{oursbg} \quad MSX, exponential only & 0.034{\scriptsize\,$\pm$0.009} & \underline{0.015}{\scriptsize\,$\pm$0.004} & 0.416{\scriptsize\,$\pm$0.008} & 11.51{\scriptsize\,$\pm$0.46} & n/a \\
\addlinespace[3pt]\multicolumn{6}{l}{\textit{S4: multiscale network}} \\[1pt]
\quad Exp.\ kernel MLE & 0.173{\scriptsize\,$\pm$0.041} & 0.095{\scriptsize\,$\pm$0.025} & 1.635{\scriptsize\,$\pm$0.040} & 711.28{\scriptsize\,$\pm$53.09} & 1.000{\scriptsize\,$\pm$0.000} \\
\quad ADM4 & 0.399{\scriptsize\,$\pm$0.003} & 0.255{\scriptsize\,$\pm$0.001} & 0.410{\scriptsize\,$\pm$0.001} & 3.13{\scriptsize\,$\pm$0.25} & 1.000{\scriptsize\,$\pm$0.000} \\
\quad Hawkes EM (nonparam.) & 0.153{\scriptsize\,$\pm$0.012} & 0.045{\scriptsize\,$\pm$0.010} & 0.358{\scriptsize\,$\pm$0.012} & 4.55{\scriptsize\,$\pm$0.29} & 1.000{\scriptsize\,$\pm$0.000} \\
\quad NPHC (cumulants) & 1.030{\scriptsize\,$\pm$0.304} & 0.211{\scriptsize\,$\pm$0.186} & -- & -- & 0.588{\scriptsize\,$\pm$0.155} \\
\quad Conditional law & 0.464{\scriptsize\,$\pm$0.286} & 0.139{\scriptsize\,$\pm$0.050} & 1.059{\scriptsize\,$\pm$0.074} & 114.94{\scriptsize\,$\pm$1.20} & 0.831{\scriptsize\,$\pm$0.087} \\
\rowcolor{oursbg} \quad MSX-auto & \textbf{0.084}{\scriptsize\,$\pm$0.014} & \underline{0.016}{\scriptsize\,$\pm$0.006} & \textbf{0.128}{\scriptsize\,$\pm$0.026} & \textbf{0.32}{\scriptsize\,$\pm$0.08} & 1.000{\scriptsize\,$\pm$0.000} \\
\rowcolor{oursbg} \quad MSX ($R=2$) & 0.093{\scriptsize\,$\pm$0.015} & 0.021{\scriptsize\,$\pm$0.005} & 0.173{\scriptsize\,$\pm$0.024} & 0.54{\scriptsize\,$\pm$0.12} & 1.000{\scriptsize\,$\pm$0.000} \\
\rowcolor{oursbg} \quad MSX ($R=4$) & 0.101{\scriptsize\,$\pm$0.019} & 0.027{\scriptsize\,$\pm$0.007} & 0.207{\scriptsize\,$\pm$0.035} & 0.70{\scriptsize\,$\pm$0.11} & \underline{1.000}{\scriptsize\,$\pm$0.000} \\
\rowcolor{oursbg} \quad MSX, exponential only & \underline{0.084}{\scriptsize\,$\pm$0.014} & \textbf{0.016}{\scriptsize\,$\pm$0.006} & \underline{0.128}{\scriptsize\,$\pm$0.026} & \underline{0.32}{\scriptsize\,$\pm$0.08} & \textbf{1.000}{\scriptsize\,$\pm$0.000} \\
\addlinespace[3pt]\multicolumn{6}{l}{\textit{S5: 10-dim.\ power-law network}} \\[1pt]
\quad Exp.\ kernel MLE & 0.337{\scriptsize\,$\pm$0.153} & 0.103{\scriptsize\,$\pm$0.032} & 1.302{\scriptsize\,$\pm$0.897} & 426.80{\scriptsize\,$\pm$368.57} & 0.977{\scriptsize\,$\pm$0.050} \\
\quad ADM4 & 0.167{\scriptsize\,$\pm$0.003} & 0.096{\scriptsize\,$\pm$0.003} & 0.319{\scriptsize\,$\pm$0.002} & 26.51{\scriptsize\,$\pm$0.29} & 1.000{\scriptsize\,$\pm$0.000} \\
\quad Hawkes EM (nonparam.) & \underline{0.085}{\scriptsize\,$\pm$0.008} & 0.045{\scriptsize\,$\pm$0.008} & 0.280{\scriptsize\,$\pm$0.010} & 5.23{\scriptsize\,$\pm$0.25} & 1.000{\scriptsize\,$\pm$0.000} \\
\quad NPHC (cumulants) & 1.348{\scriptsize\,$\pm$0.477} & 0.520{\scriptsize\,$\pm$0.372} & -- & -- & 0.672{\scriptsize\,$\pm$0.042} \\
\quad Conditional law & 0.461{\scriptsize\,$\pm$0.047} & 0.037{\scriptsize\,$\pm$0.026} & 0.641{\scriptsize\,$\pm$0.025} & 78.64{\scriptsize\,$\pm$4.79} & 0.950{\scriptsize\,$\pm$0.033} \\
\rowcolor{oursbg} \quad MSX-auto & \textbf{0.045}{\scriptsize\,$\pm$0.004} & \textbf{0.015}{\scriptsize\,$\pm$0.005} & \textbf{0.079}{\scriptsize\,$\pm$0.005} & \textbf{0.74}{\scriptsize\,$\pm$0.10} & 1.000{\scriptsize\,$\pm$0.000} \\
\rowcolor{oursbg} \quad MSX ($R=2$) & 0.111{\scriptsize\,$\pm$0.013} & 0.040{\scriptsize\,$\pm$0.013} & 0.180{\scriptsize\,$\pm$0.017} & 1.19{\scriptsize\,$\pm$0.07} & 1.000{\scriptsize\,$\pm$0.000} \\
\rowcolor{oursbg} \quad MSX ($R=4$) & 0.122{\scriptsize\,$\pm$0.008} & 0.051{\scriptsize\,$\pm$0.013} & 0.218{\scriptsize\,$\pm$0.015} & 1.63{\scriptsize\,$\pm$0.09} & \underline{1.000}{\scriptsize\,$\pm$0.000} \\
\rowcolor{oursbg} \quad MSX, exponential only & 0.099{\scriptsize\,$\pm$0.014} & \underline{0.030}{\scriptsize\,$\pm$0.014} & \underline{0.145}{\scriptsize\,$\pm$0.017} & \underline{0.82}{\scriptsize\,$\pm$0.08} & \textbf{1.000}{\scriptsize\,$\pm$0.000} \\
\bottomrule
\end{tabular}
\end{center}
